% profiles/american-university.tex — American University
% ============================================================================
% source: https://www.american.edu/provost/ogps/graduate-studies/etd/etd-style-guide.cfm
%         (Office of Graduate Studies, "Electronic Theses and Dissertations
%          Style Guide" — the filing requirements)
% source: https://www.american.edu/provost/ogps/graduate-studies/etd/guide.cfm
%         (the same guide's submission half: the step-by-step ETD process)
% source: https://www.american.edu/provost/ogps/graduate-studies/etd/upload/etd-appendix-a-11-23.pdf
%         ("APPENDIX A: SAMPLE PAGES" — the annotated specimen title page the
%          style guide points at for the degree-statement wording. Used here as
%          EVIDENCE OF LAYOUT, which is what RESEARCHING.md allows an official
%          specimen to be. Read as text AND rendered at 130dpi, because the
%          signature rules are drawn lines that no text extractor reports.)
% source: https://www.american.edu/provost/ogps/graduate-studies/etd/
%         (the ETD hub, for the template and appendix links)
% source: https://www.american.edu/digitalaccessibility/
%         (checked for an accessibility requirement; see the note below)
% accessed: 2026-09-22
% checked:  texlib-thesis-profile-round
%
% VERIFIED and set here: geometry, spacing, the title page, and the absence of
% a separate approval page.
% NOT set, and deliberately:
%   * \thesissetchapteropening — the deeper chapter opening is OPTIONAL at AU,
%     not required. See the margin note below.
%   * \thesisrequiretagging / \thesisrequiredocumentlanguage /
%     \thesisrequirepdfstandard — AU's ETD requirements name none of them.
%
% CURRENCY. Neither the style guide nor its appendices carries a revision date
% in the page body. The specimen's filename (`etd-appendix-a-11-23.pdf') reads
% as November 2023, which is the only date evidence either document offers.
% Treat the wording as current-but-undated rather than freshly dated.
%
% ONE REQUIREMENT THIS PROFILE CANNOT EXPRESS, and it is met anyway. AU: "A
% single font should be used throughout the entire thesis or dissertation,
% including text, captions, references, labels and headings", and "Fonts must
% be embedded in the document" (TrueType). texlib-thesis sets Latin Modern
% through fontspec for the whole document and embeds every face, so a document
% built with this class satisfies both. AU *recommends* 10-point Arial or
% 12-point Times New Roman; a recommendation is not a requirement, so nothing
% here changes the face. Anyone who wants the literal recommendation can add
% \setmainfont{TeX Gyre Termes} after \documentclass.
%
% SIGNATURES ARE STILL PHYSICAL, and that is the single most consequential fact
% in this file. The ETD process is: "Student collects committee members'
% signatures on a printed copy of the thesis title page", then "Student scans
% signed title page and merges the page into thesis or dissertation PDF
% document", and the submission "must ... contain[] the title page signed by
% the student's committee members and the Dean of the school or college." So
% the page this profile typesets is printed, signed by hand, scanned, and put
% back in place of the typeset one. A scanned page is an image, and an image of
% text is the least accessible thing a PDF can contain — this class's tagging
% cannot follow the document through that step. Nothing can be done about it
% from here; it is written down so nobody is surprised by it.

\thesisinstitution{American University}

% "All pages must have margins of no less than one inch on all sides (left,
% right, top and bottom)."
%
% That is a FLOOR with no stated ceiling: margins "may be wider, but not
% narrower than one inch", and must stay consistent throughout. The profile
% takes the one-inch minimum.
%
% THE CHAPTER-OPENING RULE IS A CHOICE, NOT A REQUIREMENT, and that is why
% \thesissetchapteropening is absent. AU permits a wider top margin on chapter
% title pages -- up to two inches if desired -- provided the variation is
% uniform across all major heading pages. Michigan, UNC and Georgia Tech
% *require* two inches there; AU only allows it. A flat 1in is in spec, so the
% class's default stands. A student who wants the deeper opening adds
%     \thesissetchapteropening{2in}
% to their own preamble, which is equally in spec as long as they leave it on
% for every chapter and major heading.
\thesissetgeometry{left=1in, right=1in, top=1in, bottom=1in}

% "All text must be double-spaced, with the following exceptions (these are
% optional, but may enhance the readability of your document)" -- block
% quotations over four lines, reference/footnote/endnote entries, table-of-
% contents titles, and multi-line headings, captions and titles may each be
% single-spaced, and triple-spacing before a subheading is allowed. Every
% exception is permissive, so there is nothing to encode but the rule itself.
\thesissetspacing{double}

% --- Accessibility -----------------------------------------------------------
% AU's ETD requirements say NOTHING about accessibility. The style guide, its
% submission half and its appendices were read end to end on the access date
% and contain no mention of a tagged PDF, alt text, a document language, a
% screen reader, WCAG or a PDF standard. The only file-format rules are
% ProQuest's: a single PDF, no password protection, no digital signature, no
% security settings that prevent printing, and embedded fonts.
%
% AU does publish a general "Creating Accessible PDFs" guide under
% /digitalaccessibility/, and it is genuinely about tagging, reading order and
% alt text -- but it is a how-to for staff and faculty authoring documents, it
% states no standard (no WCAG level, no Title II date), and nothing in the ETD
% requirements points at it or makes it binding on a filed thesis. Declaring
% \thesisrequiretagging on the strength of a separate how-to page would be
% putting words in the Office of Graduate Studies' mouth.
%
% So this is recorded as an absence, in the sense RESEARCHING.md distinguishes:
% the section WAS read, and it asks for nothing.
\thesisaccessibilityrequirement{Not published}
	{American University's ETD Style Guide states no accessibility requirement
	 for a filed thesis or dissertation: no tagged PDF, no alt text, no
	 document language, no WCAG level, no PDF standard. The university
	 publishes a general "Creating Accessible PDFs" guide elsewhere, but the
	 filing requirements neither cite it nor make it binding. Note also that AU
	 requires the signed title page to be scanned and merged into the submitted
	 PDF, which puts an image of text into every filed document.}

% --- Title page --------------------------------------------------------------
% AU's title page IS the approval page: the committee's and the Dean's
% signature lines are on it. The style guide fixes the content and the order --
%
%   1. "The full title of the thesis / dissertation, typed in all capital
%      letters and centered on the page."
%   2. "The word 'By' and the student's name as it appears in university
%      records, centered on the page below the title."
%   3. "A standard degree statement centered on the page, below the student's
%      name (see sample page in Appendix A for proper wording). Master's degree
%      candidates should use the word 'thesis' and doctoral candidates should
%      use the word 'dissertation.'"
%   4. "An alphabetical list of all members of the thesis/dissertation
%      committee, but with the Chair first", on the right, each name beneath a
%      signature line.
%   5. On the left, "the words 'Dean of' [student's school or college], typed
%      below a signature line and above a line for the date the page was
%      signed."
%   6. "The year (not month or day) the student's degree will be granted, and
%      the words American University, Washington D.C. 20016, centered at the
%      bottom of the page."
%
% -- and the specimen supplies the wording the guide defers on. The eleven
% centred lines below are that specimen, transcribed:
%
%       TITLE ... IN ALL CAPITAL LETTERS / DOUBLE-SPACED AND CENTERED ON PAGE
%       By
%       [Your Name]
%       Submitted to the
%       Faculty of the [Your School or College]
%       of American University
%       in Partial Fulfillment of
%       the Requirements for the Degree of
%       [Your Degree (ex. Master of Arts)]
%       In
%       [Your Discipline (as it will appear on your diploma)]
%
% The page is DOUBLE-spaced -- the specimen's own title text says so in as many
% words -- which is unusual; most schools in this directory single-space the
% title page. The signature blocks are single-spaced so each name sits directly
% under its rule, as in the specimen.
%
% TWO SMALL DISCREPANCIES, neither resolved silently:
%   * The prose says "American University, Washington D.C. 20016" on one line
%     with no comma after Washington; the specimen prints two lines, "American
%     University" then "Washington, D.C. 20016", with the comma. The specimen
%     is what is reproduced here, on the RESEARCHING.md rule that a specimen is
%     evidence of layout.
%   * The specimen places the label "Chair:" slightly to the LEFT of the
%     committee block's signature rules, an offset neither source states in
%     prose. Here it is the first line of that block. If a reviewer wants the
%     specimen's exact inset, it is a \hspace* on that one line.
%
% SO AN AMERICAN UNIVERSITY DOCUMENT DECLARES ITS COMMITTEE LIKE THIS:
%     \committeemember{Chair Name, Ph.D.}{Chair}
%     \committeemember{Second Name, Ph.D.}{}
%     \committeemember{Third Name, Ph.D.}{}
% -- chair first, the rest alphabetical, the role left empty for everyone but
% the chair, because AU labels only the chair.
%
% TWO FIELDS AU NEEDS THAT THE CLASS HAS NO SLOT FOR:
%   * the school or college, supplied here as \thesisschool. Left unset it
%     prints AU's own specimen placeholder, so an unfilled page is obviously
%     unfilled rather than quietly wrong.
%   * the year alone. AU says "The year (not month or day)", and the class's
%     \graddate is a free string whose default is "May, 20XX". At AU, set it to
%     the bare year: \graddate{2027}.
\providecommand{\thesis@school}{[Your School or College]}
\providecommand{\thesisschool}[1]{\renewcommand{\thesis@school}{#1}}
\renewcommand{\thesistitlepage}{%
	\loadgeometry{nopagenumber}%
	\begin{titlepage}%
		\thispagestyle{empty}\doublespacing\centering
		\thesis@tagtitle{Title}{\MakeUppercase{\@title}}\par
		By\par
		\@author\par
		Submitted to the\par
		Faculty of the \thesis@school\par
		of American University\par
		in Partial Fulfillment of\par
		the Requirements for the Degree of\par
		\thesis@degree\par
		In\par
		\thesis@program\par
		% Committee block: right-hand column, each name under its own signature
		% rule, the chair's role label above the first rule.
		\begingroup
		\renewcommand{\thesis@memberline}[2]{%
			\ifx\relax##2\relax\else ##2:\par\fi
			\rule{3in}{0.4pt}\par
			##1\par
			\vspace{1.2\baselineskip}%
		}%
		\hfill\begin{minipage}{3.2in}%
			\singlespacing\raggedright
			\thesis@committee
		\end{minipage}\par
		\endgroup
		% Dean's signature line, then the date line, both at the left margin.
		\begin{minipage}{\textwidth}%
			\singlespacing\raggedright
			\rule{3in}{0.4pt}\par
			Dean of the \thesis@school\par
			\vspace{2.4\baselineskip}%
			\rule{3in}{0.4pt}\par
			Date\par
		\end{minipage}\par
		\vspace{\baselineskip}
		{\centering
		\thesis@date\par
		American University\par
		Washington, D.C. 20016\par
		}%
		\vspace*{\fill}
	\end{titlepage}%
	\loadgeometry{normal}%
}

% --- Committee approval page --------------------------------------------------
% The filed document contains no separate one, because the approval is on the
% title page above. This is a stronger case for \thesisnoapprovalpage than the
% others in this directory: at Johns Hopkins and Texas A&M the neutral fallback
% would merely add a page the school does not ask for, whereas here it would
% print a SECOND set of committee signature rules in a document whose actual
% signature page is the title page. A student could sign the wrong leaf and
% have the filing rejected for it.
%
% AU also does not list an approval page anywhere in its front matter: Appendix
% A's sample pages run title page, copyright page, abstract, table of contents,
% and the sample table of contents itself lists ABSTRACT, ACKNOWLEDGMENTS, LIST
% OF TABLES, LIST OF ILLUSTRATIONS and then the chapters.
%
% As with Texas A&M, note what this does NOT rest on: AU nowhere writes the
% sentence "do not include an approval page" the way Washington does. It rests
% on where the signatures go. A reviewer who wants that exact sentence before
% accepting \thesisnoapprovalpage should downgrade this to unset -- the neutral
% page is the conservative failure.
\thesisnoapprovalpage
